\documentclass{article}

% Submission version:
\usepackage{neurips_2026}
%
% Final camera-ready version:
% \usepackage[final]{neurips_2026}
%
% Preprint version:
% \usepackage[preprint]{neurips_2026}

\usepackage[T1]{fontenc}
\usepackage[utf8]{inputenc}
\usepackage{url}
\usepackage{silence}
\WarningFilter{latex}{Command \showhyphens has changed}
\usepackage{microtype}
\usepackage{nicefrac}
\usepackage{amsfonts}
\usepackage{amsmath,amssymb,amsthm,mathtools}
\usepackage{booktabs}
\usepackage{array}
\usepackage{enumitem}
\usepackage{xcolor}
\usepackage{graphicx}
\usepackage{subcaption}
\usepackage{float}
\usepackage{hyperref}
\usepackage[nameinlink,noabbrev]{cleveref}

\hypersetup{
  colorlinks=true,
  linkcolor=blue!50!black,
  citecolor=blue!50!black,
  urlcolor=blue!50!black
}

\newcommand{\dataset}[1]{\textsc{#1}}
\newcommand{\method}{\textsc{Volume Probe}}
\newcommand{\fiber}{\textsc{Fiber Bundle Test}}
\newcommand{\dimhat}{\widehat{d}}
\newcommand{\irr}{\mathrm{irr}}
\newcommand{\R}{\mathbb{R}}
\newcommand{\vol}{\operatorname{vol}}
\IfFileExists{imgs/neurips_submission/coco_dinov3_huge_fiber_irregularity_heatmap.png}{
  \newcommand{\figroot}{imgs/neurips_submission}
}{
  \newcommand{\figroot}{../imgs/neurips_submission}
}
\newcommand{\widepaperfig}[1]{\includegraphics[width=\linewidth,height=0.82\textheight,keepaspectratio]{#1}}

\newtheorem{definition}{Definition}
\newtheorem{theorem}{Theorem}

% Avoid page-stretching warnings on float-heavy pages.
\raggedbottom

\title{Stratified Geometry in Vision Representations}
\author{
  Anonymous Authors \\
  Paper under double-blind review
}

\begin{document}
\maketitle

\begin{abstract}
We test whether dense vision representations are better described by locally stratified geometry than by one global manifold dimension. Building on the fiber-bundle volume-growth test of \citet{robinson2025token}, we analyze frozen DINOv3, SAM, SigLIP2, and AIMv2 patch-token clouds on \dataset{COCO}, with \dataset{STL10}, \dataset{DTD}, and \dataset{EuroSAT} stress tests and an auxiliary VAR-d30 generation probe. For each token cloud, we estimate local dimension, corrected slope-increase violations of the fiber-bundle null, same-image locality, and a raw-patch sparse reconstruction burden using OMP; because each token maps back to an image patch, all diagnostics can be inspected spatially. Across models and datasets, local dimension, violation rate, locality, and sparse complexity separate from one another. DINOv3 is highly image-local, SAM is especially sensitive on segmentation-like and out-of-domain structure, and SigLIP2/AIMv2 occupy higher-dimensional local regimes with different sparse-complexity profiles. The main conclusion is simple: modern dense vision spaces contain heterogeneous local structure that a single intrinsic-dimension estimate hides.
\end{abstract}

\section{Introduction}
The \emph{manifold hypothesis} says that high-dimensional data concentrate near a much lower-dimensional set. It is a useful organizing idea, but modern representation spaces often need a more local description. One global intrinsic-dimension estimate can hide neighborhoods whose scaling behavior, effective dimension, or visual content changes from point to point \citep{loaiza-ganem2024deep,levina2005maximum,facco2017estimating,pope2021intrinsic,tempczyk2022lidl,tempczyk2026benchmarks}.

We therefore treat the manifold hypothesis as a local statistical null rather than as a global slogan. Around each anchor, a regular local model predicts that empirical ball volume should grow like a stable power law in radius, with only the allowed flattening captured by the fiber-bundle formulation. The test asks a directional question: after multiple-testing correction, does a later radius window grow faster than an earlier one? A positive answer marks a local failure of the tested regular model, not a claim that the entire representation is non-manifold.

Recent work on language-token embeddings sharpens this local view: compare each token's log-volume/log-radius curve with a permissive fiber-bundle null, and flag statistically significant slope increases as violations of that null \citep{robinson2025token}. The point is not merely that tokens are discrete. It is that learned neighborhoods can fail even a smooth local model. This paper asks the analogous question for dense vision representations.

Vision gives a cleaner observational handle than a vocabulary table. A transformer or segmentation encoder produces hundreds of patch tokens per image, and each token remains tied to a visible region: an object interior, boundary, texture, occlusion, or background transition. This makes the sampling denser: one image supplies a complete grid of nearby anchors that share scene context while varying in local visual content, so a small aligned image set yields thousands of spatially interpretable local neighborhoods. We can therefore estimate local geometry in feature space and project the result back to the image. That spatial link lets us ask not only whether a representation is locally regular, but where the irregularities occur and whether the corresponding raw patches are visually simple or complex.

We study frozen DINOv3, SAM, SigLIP2, and AIMv2 checkpoints on image-aligned \dataset{COCO} subsets \citep{simeoni2025dinov3,kirillov2023segment,tschannen2025siglip2,elnouby2024aimv2,lin2014microsoft}. We add \dataset{STL10}, \dataset{DTD}, and \dataset{EuroSAT} as domain shifts, and include VAR-d30 as an auxiliary autoregressive vision-token probe \citep{tian2024visual}. The claim is deliberately local and diagnostic: dense vision representations are not well summarized by one global dimension number. Their patch-token geometry is better treated as stratified, with local scaling and reconstruction complexity varying across image regions, models, and datasets.

\paragraph{Contributions.}
\begin{itemize}[leftmargin=1.5em]
\item We adapt the fiber-bundle volume-growth test of \citet{robinson2025token} to dense vision patch-token clouds.
\item We evaluate frozen DINOv3, SAM, SigLIP2, AIMv2, and VAR-d30 spaces on \dataset{COCO}, with stress tests on \dataset{STL10}, \dataset{DTD}, and \dataset{EuroSAT}.
\item We project local dimension, corrected slope-increase violations, same-image locality, and sparse raw-patch reconstruction burden back to image patches.
\item We show that these diagnostics separate: global dimension is a poor summary, and patch-projected geometry exposes model- and region-specific stratification.
\end{itemize}

\section{Method}
\paragraph{Overview.}
The unit of analysis is a spatial patch token. For each frozen model, an image yields token vectors $z_i$ and corresponding raw RGB patches $x_i$. We use this pairing for three diagnostics: local dimension from neighbor-count growth, fiber-bundle violations from forbidden slope increases in the same volume curve, and sparse raw-patch complexity from a local OMP dictionary. These diagnostics are intentionally separate. Local dimension measures feature-space volume growth; fiber violations measure scale changes in that growth; sparse complexity measures visible variation inside the token neighborhood.

\paragraph{Geometric null.}
Let $Z=\{z_1,\ldots,z_n\}\subset \R^\ell$ be the frozen patch-token cloud. The fiber-bundle null says that a regular neighborhood may flatten as radius grows, but it should not become steeper: a later slope increase means the neighborhood has opened into directions that were not visible locally \citep{robinson2025token}. We do not estimate the reach parameter directly. We use a conservative neighborhood range and interpret corrected slope increases as evidence that this local range is not well described by the tested regular model.

\begin{definition}[Embedded patch-token sample]
For a frozen vision encoder, the embedded sample is the finite set $Z=\{z_i\}_{i=1}^n\subset \R^\ell$ of patch-token vectors. Each $z_i$ has a source image id, a patch-grid coordinate, and a raw image patch $x_i$. Geometric tests use Euclidean distances in feature space; image patches are used for projection and sparse-reconstruction diagnostics.
\end{definition}

\begin{definition}[Empirical ball volume]
For an anchor token $z_i$ and radius $r\ge0$, the empirical volume of the feature-space ball is the neighbor count
\begin{equation}
  \widehat{V}_i(r)
  =
  \left|\{j\ne i:\|z_j-z_i\|_2\le r\}\right|.
\end{equation}
Equivalently, if $r_k(i)$ is the distance from $z_i$ to its $k$th nearest neighbor, then $\widehat{V}_i(r_k(i))=k$. Sorted distances supply the radii and sorted indices supply the empirical volumes.
\end{definition}

\paragraph{Local volume scaling.}
The dimension estimate comes from the standard small-ball volume expansion used in Riemannian geometry and intrinsic-dimension estimation \citep{gray1974volume,tenenbaum2000global,levina2005maximum,facco2017estimating}. If points are sampled near a locally $d$-dimensional set, then the volume of a small ball grows like
\begin{equation}
  \vol(B_r(z_i)) = C_i r^d + O(r^{d+1}),
\end{equation}
and the empirical count $\widehat{V}_i(r)$ approximates that local volume. Taking logs gives $\log \widehat{V}_i(r) \approx c_i + d\log r$, so the slope of a log-volume/log-radius curve estimates local dimension. In practice we do not choose radii directly. The pair $(\log r_k(i), \log k)$ traces the empirical volume curve around token $i$.

\begin{definition}[Local dimension estimate]
Given neighborhood indices $k\in[k_{\min},k_{\max}]$, the local dimension estimate at $z_i$ is the slope
\begin{equation}
  \dimhat_i = \widehat{\beta}_{1,i},
  \qquad
  \log k = \beta_{0,i}+\beta_{1,i}\log r_k(i)+\varepsilon_{k,i}.
\end{equation}
We treat $\dimhat_i$ as a local diagnostic, not as a global manifold dimension. It can vary across tokens in the same image.
\end{definition}

\begin{theorem}[Fiber-bundle slope rule, after \citet{robinson2025token}]
\label{thm:fiber-monotone}
For a local neighborhood that behaves like the fiber-bundle model in \citet{robinson2025token}, the log-volume/log-radius curve may change slope as radius grows, but the slope should stay the same or decrease. A statistically significant slope increase is therefore evidence against the fiber-bundle null.
\end{theorem}

\paragraph{Fiber violations and irregularity.}
The local dimension estimate uses one slope as a summary. The fiber test asks whether that slope changes in the direction ruled out by \Cref{thm:fiber-monotone}. We record ordinary change points when the curve changes shape, but count a fiber violation only when the later radius window is steeper than the earlier one.

\begin{definition}[Patch-token irregularity score]
For each anchor token $z_i$, we scan adjacent windows on its log-volume/log-radius curve and test only the fiber-violating direction: the later window has larger slope than the earlier window. We apply Holm-Bonferroni correction \citep{holm1979simple} to these slope-increase tests, as in \citet{robinson2025token}. Let $p_i^+$ be the smallest adjusted p-value for token $z_i$. The irregularity score is
\begin{equation}
  \irr_i =
  \begin{cases}
  -\log_{10}(p_i^+), & p_i^+<\alpha,\\
  0, & p_i^+\ge \alpha.
  \end{cases}
\end{equation}
\end{definition}

Thus $\irr_i$ is zero unless a corrected slope increase is significant; larger values mean stronger evidence for the forbidden direction. The irregular ratio is the fraction of patch tokens with $\irr_i>0$. Heatmaps project this score back to image space.

\paragraph{Local sparse complexity.}
The sparse probe asks how much raw visual variation lies inside a token neighborhood. Sparse coding and dictionary learning measure whether signals can be represented with a small number of local basis elements \citep{olshausen1996emergence,mairal2009online}. We use Orthogonal Matching Pursuit (OMP) \citep{pati1993orthogonal,tropp2007signal} because it gives an interpretable code length: the number of atoms needed to hit a fixed reconstruction error.

For the fixed-neighborhood probe, $\mathcal{N}_i$ is the set of nearest token neighbors of anchor $i$. We use $k=128$ neighbors, a 128-atom local PCA dictionary, and OMP sparsity cap 64. For each anchor, we collect the corresponding raw image patches $\{x_j:j\in\mathcal{N}_i\}$, center and normalize them, build a dictionary from the leading right singular vectors of the local patch matrix, and reconstruct each patch with OMP. The required sparsity $s_{ij}$ is the smallest support size that reaches target relative residual $\tau$, subject to the cap. The token-level sparse complexity is
\begin{equation}
  S_i=\frac{1}{|\mathcal{N}_i|}\sum_{j\in\mathcal{N}_i} s_{ij},
\end{equation}
where high $S_i$ means the token neighborhood contains raw patches that are hard to explain with a very sparse local code. We report a strict baseline $\tau=0.15$ and sweep $\tau\in\{0.15,0.20,0.25,0.30,0.35,0.40\}$ to separate regional variation from cap saturation.

We also run an expanding-neighborhood probe for the spatially middle patch token. For volumes $8,16,\ldots,512$, we take the nearest token neighborhood of that size, fit a fresh local dictionary, and rerun OMP. If required sparsity still rises, the effect is not caused by forcing all scales to share one global dictionary.

\paragraph{Reading the visualizations.}
The visualizations are part of the evidence. Fiber heatmaps project $\irr_i$ back to the source image grid. Dimension galleries show the raw patches selected as low-, middle-, high-dimension, or high-violation examples. Sparse heatmaps project $S_i$ back to image space, and middle-patch volume curves show how OMP sparsity changes as neighborhoods grow. Together they keep local feature-space scaling, corrected slope changes, and raw-patch reconstruction complexity distinct.

\section{Experimental Setup}
\paragraph{Data.}
The main experiments use \dataset{COCO} images resized to $224\times224$. The goal is not to train or rank classifiers, but to use dense image tokens as a medium for local geometric diagnostics. Patch-size-16 encoders give a $14\times14$ grid, or 196 tokens per image. AIMv2 and VAR use $16\times16$ analysis grids, or 256 tokens per image. Most \dataset{COCO} comparisons use the same 16 image-aligned examples; the DINOv3 ViT-L/16 row uses a preceding 21-image run as a scale comparison.

We run the same frozen-encoder pipeline on \dataset{STL10}, \dataset{DTD}, and \dataset{EuroSAT}. These are stress tests: \dataset{STL10} gives object-centered images, \dataset{DTD} emphasizes textures, and \dataset{EuroSAT} gives satellite scenes. SAM runs outside \dataset{COCO} use frozen image-encoder patch tokens only, without box prompts or mask previews.

\paragraph{Models for Evaluation.}
We use ViT-style checkpoints because they expose one vector per image patch \citep{dosovitskiy2021image}. DINOv3 provides strong self-supervised visual features \citep{simeoni2025dinov3}; SAM provides segmentation-oriented encoder tokens and COCO mask previews for context \citep{kirillov2023segment}; SigLIP2 and AIMv2 test modern multimodal dense encoders \citep{tschannen2025siglip2,elnouby2024aimv2}; VAR-d30 is an auxiliary autoregressive visual-token generator rather than a matched encoder baseline \citep{tian2024visual}. All geometry and sparse probes use frozen representations. Table~\ref{tab:checkpoint-settings} gives the exact checkpoints and settings.

\begin{table}[H]
\centering
\scriptsize
\caption{Frozen checkpoints and settings used in the COCO experiments. All runs use pretrained encoders and set training epochs to zero for the representation probes.}
\label{tab:checkpoint-settings}
\setlength{\tabcolsep}{3pt}
\begin{tabular}{@{}p{0.13\linewidth}p{0.31\linewidth}p{0.18\linewidth}p{0.30\linewidth}@{}}
\toprule
Run & Checkpoint & Token sample & Main settings \\
\midrule
DINOv3-H+ &
\texttt{vit\_huge\_plus\_patch16\_}\newline
\texttt{dinov3.lvd1689m} from \texttt{timm} &
16 images, 3136 tokens, 1280-d embeddings &
Patch size 16; image size 224; \texttt{vol\_min=8}, \texttt{vol\_max=64}, \texttt{ws=4}, \texttt{alpha=0.005}; high-capacity residual sweep uses fixed $k=128$, 128 local PCA atoms, and OMP cap 64. \\
DINOv3-L &
\texttt{vit\_large\_patch16\_}\newline
\texttt{dinov3.lvd1689m} from \texttt{timm} &
21 images, 4116 tokens, 1024-d embeddings &
Scale-comparison run with the same DINO probe settings and local OMP dictionary configuration. \\
SAM-H &
\texttt{facebook/sam-vit-huge} from Hugging Face &
16 images, 3136 analysis tokens &
Patch-size-16 analysis grid; \texttt{vol\_min=8}, \texttt{vol\_max=64}, \texttt{ws=8}, \texttt{alpha=0.005}; high-capacity residual sweep uses fixed $k=128$, 128 local PCA atoms, and OMP cap 64; COCO box prompts capped at 16 per image with mask threshold 0.25 for mask previews. \\
SigLIP2-B &
\texttt{google/siglip2-base-patch16-224} from Hugging Face &
16 images, 3136 tokens, 768-d embeddings &
Patch size 16; image size 224; \texttt{vol\_min=8}, \texttt{vol\_max=64}, \texttt{ws=4}, \texttt{alpha=0.005}; sparse probe uses fixed $k=128$, 128 local PCA atoms, residual target 0.30, and OMP cap 64. \\
AIMv2-L &
\texttt{apple/aimv2-large-patch14-224} from Hugging Face &
16 images, 4096 tokens, 1024-d embeddings &
Patch size 14; $16\times16$ analysis grid; \texttt{vol\_min=8}, \texttt{vol\_max=64}, \texttt{ws=4}, \texttt{alpha=0.005}; sparse probe uses fixed $k=128$, 128 local PCA atoms, residual target 0.30, and OMP cap 64. \\
VAR-d30 &
Autoregressive visual-token checkpoint &
16 images, 4096 tokens &
Auxiliary generation-token probe on a $16\times16$ visual-token grid; geometry and sparse probes are reported separately from the encoder comparisons. \\
\bottomrule
\end{tabular}
\end{table}

\paragraph{Probe settings.}
The fiber probe uses image-aligned token collection, so every selected image contributes a complete patch grid. Same-image locality uses $k=16$ nearest neighbors. The sparse-complexity sweep uses fixed $k=128$ token neighborhoods, a 128-atom local PCA dictionary, and OMP cap 64 for DINOv3-H+, SAM-H, SigLIP2-B, and AIMv2-L. Each token gets its own local dictionary; the expanding-volume probe also refits the dictionary at each neighborhood size.

\section{Results}
\subsection{Summary Statistics}
\begin{table}[H]
\centering
\scriptsize
\caption{\dataset{COCO} pretrained feature probes. ``Change'' counts any significant slope change; ``fiber viol.'' counts only significant slope increases. Sparse complexity reports the high-capacity fixed-$k=128$, $\tau=0.30$ OMP probe with 128 local PCA atoms and sparsity cap 64 when available. The DINOv3-L row is retained as a scale comparison; VAR-d30 is an auxiliary visual-token generation probe.}
\begin{tabular}{@{}lrrrrrrr@{}}
\toprule
Model & Images & Tokens & Mean/med. dim & Change & Fiber viol. & Same-image & Mean $S_i$ \\
\midrule
DINOv3 ViT-L/16 & 21 & 4116 & 5.95/5.31 & 0.217 & 0.090 & 0.985 & 15.53 \\
DINOv3 ViT-H+/16 & 16 & 3136 & 6.61/5.98 & 0.200 & 0.076 & 0.979 & 14.83 \\
SAM ViT-Huge encoder & 16 & 3136 & 6.24/6.01 & 0.243 & 0.059 & 0.659 & 20.03 \\
SigLIP2-B & 16 & 3136 & 11.17/9.30 & 0.233 & 0.100 & 0.806 & 14.30 \\
AIMv2-L & 16 & 4096 & 12.30/10.29 & 0.217 & 0.102 & 0.746 & 12.64 \\
VAR-d30 & 16 & 4096 & 36.59/36.38 & 0.168 & 0.042 & 0.143 & 13.86 \\
\bottomrule
\end{tabular}
\label{tab:coco-vit-sam}
\end{table}

The first result is that a single dimension number is too coarse. DINOv3-H+ and SAM-H have similar mean/median local dimensions, 6.61/5.98 and 6.24/6.01, but different locality and change-point structure. SigLIP2-B and AIMv2-L sit in a higher-dimensional regime and have fiber-violation ratios near 0.10. VAR-d30 is higher-dimensional still, yet has the lowest fiber-violation ratio. Local dimension, generic change points, and corrected slope-increase violations are therefore separate signals.

\subsection{Neighborhood Locality}
\begin{table}[H]
\centering
\scriptsize
\caption{Additional \dataset{COCO} local-geometry diagnostics. Same-image statistics use $k=16$ nearest neighbors in feature space. ``Chart'' is the within-image local-chart score, ``multi'' is the fraction of anchors whose scaling curve has more than one detected regime, and the last two columns summarize image-to-image and within-image variation of the local dimension field.}
\begin{tabular}{@{}lrrrrrrrrr@{}}
\toprule
Model & Top-1 same & Mean same & Lift & Chart & Multi & Mean irr. & Max irr. & Image CV & Within std \\
\midrule
DINOv3 ViT-L/16 & 1.000 & 0.985 & 20.8 & 0.413 & 0.126 & 0.318 & 6.99 & 0.170 & 2.79 \\
DINOv3 ViT-H+/16 & 0.998 & 0.979 & 15.7 & 0.377 & 0.113 & 0.270 & 6.63 & 0.191 & 3.12 \\
SAM ViT-Huge encoder & 0.849 & 0.659 & 10.6 & 0.327 & 0.165 & 0.240 & 9.22 & 0.217 & 3.02 \\
SigLIP2-B & 0.951 & 0.806 & 13.0 & 0.220 & 0.144 & 0.356 & 7.88 & 0.230 & 8.64 \\
AIMv2-L & 0.915 & 0.746 & 12.0 & 0.156 & 0.130 & 0.359 & 7.28 & 0.167 & 8.75 \\
VAR-d30 & 0.225 & 0.143 & 2.3 & 0.145 & 0.096 & 0.152 & 8.00 & 0.070 & 8.46 \\
\bottomrule
\end{tabular}
\label{tab:coco-vit-sam-local}
\end{table}

Neighborhood locality shows what the volume-growth test is sampling. DINOv3 is extremely image-local, with DINOv3-L/H+ mean same-image neighbor fractions of 0.985 and 0.979. SAM-H is less image-local (0.659) but still 10.6$\times$ above chance. SigLIP2-B and AIMv2-L remain far above chance while mixing more than DINOv3. VAR-d30 is qualitatively different, with same-image fraction 0.143, consistent with a generation-token space rather than a frozen encoder neighborhood.

Same-image locality does not say whether token distances preserve raw-patch distances inside an image. We therefore also compare the pairwise distance matrix of raw RGB patches with the pairwise distance matrix of the corresponding model tokens; details and visualizations are in \Cref{app:patch-token-distance}. This diagnostic asks how much the learned feature metric warps the image's original patch geometry. A high score means that patches near an anchor in raw RGB space also tend to be near it in token space. A low score does not mean failure: it can mean that the representation has reorganized patches according to semantics, context, segmentation cues, or invariances that are not visible in raw pixel distance. In the current \dataset{COCO} runs, SAM-H preserves more within-image raw-patch distance structure than DINOv3-H+, while SigLIP2-B and AIMv2-L move furthest away from raw RGB patch geometry.

The next plot groups are a visual audit. Irregularity heatmaps show where the fiber-bundle slope rule fails after correction. Dimension galleries show the raw patches selected by the estimator. Dimension heatmaps show whether local dimension is spatially coherent. These are different quantities: high dimension means fast first-regime neighbor growth, while high irregularity means the volume curve bends in the forbidden direction.

\begin{figure}[p]
\centering
\widepaperfig{\figroot/coco_dinov3_huge_fiber_irregularity_heatmap}
\caption{\dataset{COCO} fiber-bundle violations projected back to image patches for DINOv3 ViT-H+/16. Each square is one patch token placed at its original image location. Bright overlays mark tokens whose local volume curve becomes significantly steeper at a larger radius, after Holm-Bonferroni correction; dark regions either have no detected slope increase or do not survive correction.}
\label{fig:coco-fiber-heatmaps}
\end{figure}

\begin{figure}[p]
\ContinuedFloat
\centering
\widepaperfig{\figroot/coco_sam_fiber_irregularity_heatmap}
\caption{\dataset{COCO} fiber-bundle violations projected back to image patches for SAM ViT-Huge. This is the same diagnostic as the preceding DINOv3 panel: it is not a saliency map or a class-confidence map, but a map of where the representation neighborhood stops behaving like the tested fiber-bundle null.}
\end{figure}

\begin{figure}[p]
\centering
\widepaperfig{\figroot/coco_dinov3_huge_dimension_patch_gallery}
\caption{Patch crops grouped by low, middle, high first-regime local dimension and highest slope-increase violation score for DINOv3 ViT-H+/16. Each crop is the raw image patch corresponding to a token selected by the geometric estimator.}
\label{fig:coco-dimension-patches}
\end{figure}

\begin{figure}[p]
\ContinuedFloat
\centering
\widepaperfig{\figroot/coco_sam_dimension_patch_gallery}
\caption{Patch crops grouped by low, middle, high first-regime local dimension and highest slope-increase violation score for SAM ViT-Huge. Low-dimensional crops are often visually simple or locally homogeneous, while high-dimensional crops often include boundaries, texture mixtures, or partial objects. The irregularity row is selected by corrected slope-increase evidence, not by first-regime dimension alone.}
\end{figure}

\begin{figure}[p]
\centering
\widepaperfig{\figroot/coco_dinov3_huge_fiber_dimension_heatmap}
\caption{Local dimension field for DINOv3 ViT-H+/16. Each colored square is one patch token, with color set by the first-regime log-volume/log-radius slope. DINOv3-H+ shows strong within-image coherence: many neighboring patches share similar estimated dimension, consistent with its very high same-image nearest-neighbor rate.}
\label{fig:coco-dimension-heatmaps}
\end{figure}

\begin{figure}[p]
\ContinuedFloat
\centering
\widepaperfig{\figroot/coco_sam_fiber_dimension_heatmap}
\caption{Local dimension field for SAM ViT-Huge. SAM-H has a visually more mixed field, matching the quantitative result that it crosses image boundaries earlier while still preserving object- and boundary-relevant locality. These heatmaps should be read together with \Cref{fig:coco-fiber-heatmaps}: high dimension is not the same as high irregularity.}
\end{figure}

\Cref{fig:coco-fiber-heatmaps,fig:coco-dimension-patches,fig:coco-dimension-heatmaps} show that geometry is not uniform inside an image: nearby regions can have different local dimension and different violation evidence. DINOv3 appears more image-local and spatially coherent, while SAM shows more cross-image mixing and patch-to-patch variation, consistent with a segmentation representation that groups recurring boundaries and foreground/background transitions.

\subsection{Volume Curves and Sparse Complexity}
\begin{table}[H]
\centering
\scriptsize
\caption{Expanding-neighborhood OMP probe for the image-aligned DINOv3-H+/SAM runs. For the spatially middle patch token, we increase neighborhood volume and fit a fresh local raw-patch dictionary at each volume. ``Images'' counts how many source images enter the token neighborhood, $\bar S$ is the mean required OMP sparsity, and ``hit'' is the fraction of patches reaching residual threshold $\tau=0.15$ before the sparsity cap.}
\begin{tabular}{@{}rrrrrrr@{}}
\toprule
Volume & DINO images & DINO $\bar S$ & DINO hit & SAM images & SAM $\bar S$ & SAM hit \\
\midrule
8 & 1 & 4.62 & 1.00 & 1 & 4.50 & 1.00 \\
16 & 1 & 8.75 & 1.00 & 2 & 7.12 & 1.00 \\
32 & 1 & 14.91 & 1.00 & 2 & 9.66 & 1.00 \\
64 & 1 & 22.22 & 0.83 & 3 & 14.47 & 0.25 \\
128 & 3 & 22.89 & 0.15 & 5 & 15.89 & 0.10 \\
256 & 4 & 21.54 & 0.20 & 9 & 15.11 & 0.16 \\
512 & 8 & 19.42 & 0.28 & 10 & 14.65 & 0.18 \\
\bottomrule
\end{tabular}
\label{tab:coco-volume-sparse}
\end{table}

The sparse probe asks whether feature-space neighborhoods also contain simple raw visual neighborhoods. We use the higher-capacity $k=128$ setting because the initial $k=32$ probe was too capacity-limited for SAM. If neighborhoods remain hard under this setting, they are not hard merely because the dictionary is tiny.

The high-capacity sweep confirms that sparse complexity is not just cap saturation. At $\tau=0.15$, DINOv3-H+ requires 34.67 atoms on average and SAM-H requires 42.05. At $\tau=0.30$, both models avoid cap hits while retaining regional variation: 14.83 atoms for DINOv3-H+ and 20.03 for SAM-H. We use $\tau=0.30$ for the main heatmaps.

\begin{table}[H]
\centering
\scriptsize
\caption{High-capacity residual-threshold sweep for fixed-$k=128$ local OMP probes with 128 local PCA atoms and sparsity cap 64. The dictionary and neighborhoods are held fixed; only the target relative residual $\tau$ changes. ``Patch cap'' is the fraction of individual patch reconstructions that exhaust the allowed sparsity cap.}
\begin{tabular}{@{}rccrccr@{}}
\toprule
$\tau$ & DINO $\bar S$ & DINO q10--q90 & DINO patch cap & SAM $\bar S$ & SAM q10--q90 & SAM patch cap \\
\midrule
0.15 & 34.67 & 24.61--45.35 & 0.2\% & 42.05 & 30.27--52.51 & 3.1\% \\
0.20 & 25.75 & 17.04--34.94 & 0.0\% & 32.65 & 21.99--42.21 & 0.3\% \\
0.25 & 19.41 & 12.20--27.10 & 0.0\% & 25.50 & 16.20--34.04 & 0.0\% \\
0.30 & 14.83 & 9.08--21.03 & 0.0\% & 20.03 & 11.97--27.46 & 0.0\% \\
0.35 & 11.49 & 6.99--16.44 & 0.0\% & 15.80 & 9.04--22.13 & 0.0\% \\
0.40 & 9.01 & 5.47--12.91 & 0.0\% & 12.51 & 7.00--17.78 & 0.0\% \\
\bottomrule
\end{tabular}
\label{tab:sparse-residual-sweep}
\end{table}

The expanding-neighborhood probe stress-tests raw-patch complexity. DINOv3-H+ stays within one source image up to volume 64, while its required sparsity rises from $\bar S=4.62$ to $\bar S=22.22$. SAM-H crosses images earlier, but rises more slowly, from $\bar S=4.50$ to $\bar S=14.47$. At larger volumes both models hit the dictionary budget, so the tail is a capacity warning rather than a precise estimate.

\begin{figure}[p]
\centering
\widepaperfig{\figroot/coco_dinov3_huge_fiber_volume_curve_gallery}
\caption{Representative DINOv3 ViT-H+/16 fiber-bundle volume curves projected back to their source image patches. Each row pairs a raw patch crop with its log neighbor-count versus log-radius curve. A nearly straight curve means one local dimension is a reasonable summary over the plotted radius range; a curve that steepens after a detected change point is the fiber violation.}
\label{fig:coco-volume-curves}
\end{figure}

\begin{figure}[p]
\ContinuedFloat
\centering
\widepaperfig{\figroot/coco_sam_fiber_volume_curve_gallery}
\caption{Representative SAM ViT-Huge fiber-bundle volume curves projected back to source image patches. The same slope rule is used as in the DINOv3 panel: flattening is allowed by the tested fiber-bundle null, while significant steepening indicates that the token suddenly starts seeing additional directions in feature space.}
\end{figure}

\begin{figure}[p]
\centering
\widepaperfig{\figroot/coco_sparse_residual_sweep_hicap_comparison}
\caption{Residual-threshold sensitivity for the high-capacity fixed-$k=128$ local sparse dictionary probe. DINOv3-H+ and SAM-H both relax smoothly as the residual target is loosened. The important difference from the original low-capacity setting is that SAM no longer saturates: the 64-atom cap is rarely hit even at the strict $\tau=0.15$ target.}
\label{fig:coco-sparse-residual-sweep}
\end{figure}

\begin{figure}[p]
\centering
\widepaperfig{\figroot/coco_dinov3_huge_hicap_tau_0p30_sparse_dictionary_heatmaps}
\caption{DINOv3 ViT-H+/16 patch-projected mean OMP sparsity for fixed $k=128$ neighborhoods at $\tau=0.30$, using one 128-atom local PCA dictionary per anchor and sparsity cap 64. Brighter patches need more dictionary atoms to reach the same residual target, so their feature-space neighborhood contains more diverse raw visual content.}
\label{fig:coco-sparse}
\end{figure}

\begin{figure}[p]
\ContinuedFloat
\centering
\widepaperfig{\figroot/coco_sam_hicap_tau_0p30_sparse_dictionary_heatmaps}
\caption{SAM ViT-Huge patch-projected mean OMP sparsity for fixed $k=128$ neighborhoods at $\tau=0.30$, using the same 128-atom local PCA dictionary and cap-64 protocol as the preceding DINOv3 panel. The high-capacity setting is especially important for SAM because the original low-capacity setting hid regional variation behind cap saturation.}
\end{figure}

\begin{figure}[p]
\ContinuedFloat
\centering
\includegraphics[width=\linewidth]{\figroot/coco_dinov3_huge_local_dictionary_volume_curve}
\caption{DINOv3 ViT-H+/16 middle-patch sparse-complexity volume curve. Neighborhood volume is increased while fitting a fresh local dictionary at each volume. Rising curves mean the neighborhood is becoming harder to summarize visually as more feature neighbors are included.}
\end{figure}

\begin{figure}[p]
\ContinuedFloat
\centering
\includegraphics[width=\linewidth]{\figroot/coco_sam_local_dictionary_volume_curve}
\caption{SAM ViT-Huge middle-patch sparse-complexity volume curve, using the same fresh-dictionary-per-volume protocol as the DINOv3 panel.}
\end{figure}

\begin{figure}[p]
\centering
\includegraphics[width=\linewidth]{\figroot/coco_dinov3_huge_hicap_tau_0p30_sparse_dictionary_probe}
\caption{DINOv3 ViT-H+/16 high-capacity sparse-probe dashboard at $\tau=0.30$. The dashboard summarizes the full fixed-$k=128$ sparse coding field behind \Cref{fig:coco-sparse}. DINOv3-H+ evaluates all 3136 tokens with a 128-atom local dictionary and a 64-atom cap, giving mean required sparsity 14.83 and 10--90\% range 9.08--21.03.}
\label{fig:coco-sparse-dashboards}
\end{figure}

\begin{figure}[p]
\ContinuedFloat
\centering
\includegraphics[width=\linewidth]{\figroot/coco_sam_hicap_tau_0p30_sparse_dictionary_probe}
\caption{SAM ViT-Huge high-capacity sparse-probe dashboard at $\tau=0.30$. SAM-H evaluates all 3136 tokens with a 128-atom local dictionary and a 64-atom cap, giving mean required sparsity 20.03 and 10--90\% range 11.97--27.46. Sparse reconstruction burden is a separate local signal rather than a disguised dimension or irregularity score.}
\end{figure}

Sparse plots bridge feature geometry and visible image content. High local dimension means neighbors grow quickly in feature space; high sparse complexity means the corresponding raw patches require many atoms to reconstruct. The two can coincide, but need not.

\clearpage

\subsection{Additional Vision Encoders}
SigLIP2-B and AIMv2-L extend the pipeline beyond the DINOv3/SAM contrast. Both preserve strong same-image locality, but less than DINOv3. Both also have higher local dimensions than DINOv3/SAM and the highest fiber-violation ratios in the encoder set. This is not a model-quality ranking; it shows that the diagnostic is sensitive to different representational organizations.

The sparse residual sweep adds the counterpoint: at $\tau=0.30$, SigLIP2-B requires 14.30 atoms on average and AIMv2-L requires 12.64. These values are comparable to DINOv3-H+ and below SAM-H, despite higher local dimensions and fiber-violation ratios.

\begin{figure}[p]
\centering
\widepaperfig{\figroot/coco_vision_expansion_residual_sweep_comparison}
\caption{Residual-threshold sensitivity for the additional SigLIP2-B and AIMv2-L fixed-$k=128$ local sparse dictionary probes. Both models use 128 local PCA atoms and OMP cap 64. The smooth decline in required atoms as $\tau$ is relaxed shows that the selected $\tau=0.30$ visualizations are not dominated by cap saturation.}
\label{fig:coco-vision-expansion-residual-sweep}
\end{figure}

\begin{figure}[p]
\centering
\widepaperfig{\figroot/coco_siglip2_base_fiber_irregularity_heatmap}
\caption{SigLIP2-B fiber-bundle violations projected back to the 16 \dataset{COCO} images. SigLIP2 has higher mean local dimension than DINOv3/SAM and a 0.100 fiber-violation ratio, making it a useful multimodal-encoder stress test for the same stratified-geometry pipeline.}
\label{fig:coco-siglip-aim-geometry}
\end{figure}

\begin{figure}[p]
\ContinuedFloat
\centering
\widepaperfig{\figroot/coco_aimv2_large_fiber_irregularity_heatmap}
\caption{AIMv2-L fiber-bundle violations projected back to the 16 \dataset{COCO} images. AIMv2 uses a $16\times16$ analysis grid, so it contributes 4096 tokens rather than 3136.}
\end{figure}

\begin{figure}[p]
\centering
\widepaperfig{\figroot/coco_siglip2_base_sparse_dictionary_heatmaps}
\caption{SigLIP2-B patch-projected mean OMP sparsity for fixed $k=128$ neighborhoods at $\tau=0.30$. The local dictionary protocol is the same as the DINOv3/SAM high-capacity probe.}
\label{fig:coco-siglip-aim-sparse}
\end{figure}

\begin{figure}[p]
\ContinuedFloat
\centering
\widepaperfig{\figroot/coco_aimv2_large_sparse_dictionary_heatmaps}
\caption{AIMv2-L patch-projected mean OMP sparsity for fixed $k=128$ neighborhoods at $\tau=0.30$. AIMv2 has high estimated local dimension, but its local raw-patch sparse complexity is lower than SAM-H under this OMP protocol.}
\end{figure}

\clearpage

\subsection{Cross-Dataset Stress Tests}
\begin{table}[H]
\centering
\scriptsize
\setlength{\tabcolsep}{2.5pt}
\caption{Cross-dataset frozen-encoder stress tests on \dataset{STL10}, \dataset{DTD}, and \dataset{EuroSAT}. Each row uses the same patch-token geometry pipeline as the \dataset{COCO} runs. ``Same'' is the mean same-image nearest-neighbor fraction at $k=16$. Mean $S_i$ is the fixed-neighborhood OMP sparse-complexity score from the local dictionary probe. SAM rows on these datasets use frozen SAM encoder tokens only; no box prompts or mask previews are available.}
\label{tab:cross-dataset-stress}
\begin{tabular}{@{}llrrrrrr@{}}
\toprule
Dataset & Model & Tokens & Mean/med. dim & Change & Fiber viol. & Same & Mean $S_i$ \\
\midrule
\dataset{STL10} & DINOv3-H+ & 3136 & 6.44/5.73 & 0.203 & 0.069 & 0.994 & 11.03 \\
\dataset{STL10} & SAM-H & 3136 & 5.46/4.89 & 0.278 & 0.090 & 0.704 & 13.50 \\
\dataset{STL10} & SigLIP2-B & 3136 & 7.94/6.49 & 0.207 & 0.081 & 0.834 & 10.05 \\
\dataset{STL10} & AIMv2-L & 4096 & 9.00/7.86 & 0.211 & 0.086 & 0.739 & 7.59 \\
\dataset{DTD} & DINOv3-H+ & 3136 & 6.39/6.09 & 0.169 & 0.067 & 0.858 & 6.29 \\
\dataset{DTD} & SAM-H & 3136 & 6.09/5.75 & 0.235 & 0.037 & 0.829 & 22.92 \\
\dataset{DTD} & SigLIP2-B & 3136 & 9.21/8.59 & 0.179 & 0.074 & 0.732 & 7.24 \\
\dataset{DTD} & AIMv2-L & 4096 & 9.58/9.15 & 0.186 & 0.069 & 0.620 & 6.55 \\
\dataset{EuroSAT} & DINOv3-H+ & 3136 & 8.45/7.55 & 0.173 & 0.053 & 0.974 & 16.17 \\
\dataset{EuroSAT} & SAM-H & 3136 & 4.03/3.00 & 0.425 & 0.172 & 0.841 & 27.85 \\
\dataset{EuroSAT} & SigLIP2-B & 3136 & 10.87/10.26 & 0.175 & 0.058 & 0.414 & 4.94 \\
\dataset{EuroSAT} & AIMv2-L & 4096 & 12.55/11.75 & 0.171 & 0.056 & 0.327 & 3.49 \\
\bottomrule
\end{tabular}
\end{table}

The cross-dataset results keep the same story but change the details. On \dataset{STL10}, DINOv3-H+ remains almost perfectly image-local, while SAM-H has the highest change-point and fiber-violation ratios. On \dataset{DTD}, texture neighborhoods are sparse-simple for DINOv3, SigLIP2, and AIMv2, but not for SAM. On \dataset{EuroSAT}, SAM-H is the clearest stress case: low median local dimension, high change rate, high corrected fiber-violation ratio, and the largest sparse burden. The useful point is not a single ranking; local dimension, violations, locality, and sparse complexity respond differently to domain shift.

\clearpage

\subsection{SAM Mask Context}
\begin{figure}[H]
\centering
\widepaperfig{\figroot/coco_sam_box_masks}
\caption{SAM ViT-Huge box-prompt mask previews on \dataset{COCO}. The $4\times4$ grid follows the same 16 image-aligned examples used for the token-level probes. The overlays show COCO box prompts, thresholded SAM masks when available, and the analysis patch grid. The masks are not used as labels for the fiber estimator; they only provide object-aligned context for interpreting where SAM's encoder-token geometry may be responding to object interiors, contours, or background transitions.}
\label{fig:coco-sam-masks}
\end{figure}

SAM turns the dense-vision framing into an object-aligned one. \dataset{COCO} boxes generate mask previews from the same model whose encoder tokens we analyze. The statistics still come from encoder-token neighborhoods, not mask quality, but the masks help relate geometric structure to object interiors, boundaries, and background.

\clearpage

\subsection{Patch-Size Sensitivity}
Patch size has two meanings in these experiments. The ViT/SAM architectural patch size is fixed by the pretrained checkpoint, while the \emph{analysis patch size} controls the spatial grid on which encoder features are pooled and projected back to the image. We therefore ran a small geometry-only SAM ViT-Huge sensitivity sweep over analysis patch sizes 8, 16, and 32 pixels on the same 16 \dataset{COCO} images. Sparse OMP and mask previews were disabled for this sweep, so this patch-size table has no dictionary-size hyperparameter and isolates the geometry of the pooled token grid.

\begin{table}[H]
\centering
\scriptsize
\caption{SAM ViT-Huge patch-size sensitivity on the same 16 \dataset{COCO} images. Patch size here means analysis-grid patch size after pooling the frozen SAM image-encoder feature map, not a different SAM checkpoint.}
\begin{tabular}{@{}rrrrrrrr@{}}
\toprule
Patch & Grid & Tokens & Mean/med. dim & Change & Fiber viol. & Same-image & Local chart \\
\midrule
8  & $28\times28$ & 12544 & 9.47/8.78 & 0.266 & 0.090 & 0.727 & 0.306 \\
16 & $14\times14$ & 3136  & 6.24/6.01 & 0.243 & 0.059 & 0.659 & 0.327 \\
32 & $7\times7$   & 784   & 3.88/3.97 & 0.319 & 0.098 & 0.547 & 0.443 \\
\bottomrule
\end{tabular}
\label{tab:sam-patch-size-sweep}
\end{table}

The trend is useful but not benchmark-level. Finer 8-pixel patches produce more tokens and higher estimated local dimension; coarser 32-pixel patches reduce local dimension but increase detected slope changes and fiber violations. The signal is therefore not an artifact of exactly one $14\times14$ grid.

\clearpage

\subsection{Auxiliary Autoregressive Generation Probe}
We include VAR-d30 as an auxiliary \dataset{COCO} probe, not as a matched encoder baseline \citep{tian2024visual}. It operates on a $16\times16$ grid of discrete visual codes and asks a narrower question: whether fiber irregularity in a generative token space lines up with generation uncertainty or branch diversity.

\begin{table}[H]
\centering
\scriptsize
\caption{Auxiliary VAR-d30 \dataset{COCO} generation probe from W\&B run \texttt{skt3dcv6}. The model emits discrete visual-token predictions. Entropy, negative log-likelihood (NLL), top-1 probability, and margin are measured from the next-token distribution; fiber statistics are computed on the same token representations as in the other runs.}
\begin{tabular}{@{}lrrrrrrr@{}}
\toprule
Model & Images & Tokens & Grid & Mean/med. dim & Change & Fiber viol. & Mean $S_i$ \\
\midrule
VAR-d30 generator & 16 & 4096 & $16\times16$ & 36.59/36.38 & 0.168 & 0.042 & 13.86 \\
\bottomrule
\end{tabular}
\label{tab:coco-var}
\end{table}

VAR-d30 looks different from the encoder spaces. Its same-image neighbor fraction is only 0.143, far below the frozen encoders. Its local dimension is much higher, with mean/median 36.59/36.38, but its fiber-violation ratio is lower at 0.042. High local dimension therefore does not imply many slope-increase violations.

Under teacher forcing, corrected fiber violations do not predict higher generation uncertainty. Fiber-violating tokens have slightly lower normalized entropy than non-violating tokens ($-0.017$, permutation $p=0.016$), and the high-irregularity decile does not transfer cleanly to NLL. Local dimension is the stronger generation-side signal: the high-dimension decile has lower entropy ($-0.091$, $p<10^{-4}$) and lower NLL ($-0.873$, $p<10^{-4}$). After rank-residualizing on irregularity, image id, row, and column, dimension retains larger partial Spearman correlations than irregularity.

\begin{table}[H]
\centering
\scriptsize
\caption{Aggregate teacher-forced generation diagnostics for VAR-d30 on \dataset{COCO}. Differences are first group minus second group. P-values are two-sided permutation tests over token-level metrics. Entropy is normalized by $\log$ vocabulary size.}
\begin{tabular}{@{}lrrrr@{}}
\toprule
Comparison & $\Delta$ entropy & $p$ & $\Delta$ NLL & $p$ \\
\midrule
Fiber violation vs. no violation & -0.017 & 0.016 & -0.215 & 0.167 \\
High vs. low irregularity decile & -0.010 & 0.106 & +0.085 & 0.542 \\
High vs. low dimension decile & -0.091 & $<10^{-4}$ & -0.873 & $<10^{-4}$ \\
\bottomrule
\end{tabular}
\label{tab:coco-var-generation}
\end{table}

The VAR plots are therefore a negative control for a too-simple singularity-to-polysemy story. Dimension, irregularity, and generation uncertainty do not light up in the same places: dimension is broadly high and predictive of uncertainty, irregularity is sparse, and fiber violations alone do not robustly predict local generation ambiguity.

\begin{figure}[p]
\centering
\widepaperfig{\figroot/coco_var_d30_fiber_dimension_heatmap}
\caption{VAR-d30 local dimension heatmap from W\&B run \texttt{skt3dcv6}. Many tokens live in a high-dimensional first-regime neighborhood; this measures fast local neighbor growth and should not be read as the same quantity as fiber irregularity.}
\label{fig:coco-var-geometry}
\end{figure}
\clearpage

\begin{figure}[p]
\ContinuedFloat
\centering
\widepaperfig{\figroot/coco_var_d30_fiber_irregularity_heatmap}
\caption{VAR-d30 fiber-irregularity heatmap. Only a smaller subset of tokens has corrected slope-increase evidence, meaning the local volume curve changes in the forbidden direction across radii.}
\end{figure}
\clearpage

\begin{figure}[p]
\ContinuedFloat
\centering
\widepaperfig{\figroot/coco_var_d30_dimension_patch_gallery}
\caption{VAR-d30 patch crops selected by the local dimension and irregularity diagnostics. Showing this gallery at full page width makes the token-level statistics inspectable as visible image content rather than only as numeric scores.}
\end{figure}
\clearpage

\begin{figure}[p]
\ContinuedFloat
\centering
\widepaperfig{\figroot/coco_var_d30_sparse_dictionary_probe}
\caption{VAR-d30 sparse-probe dashboard. The fixed-$k$ sparse-coding view is a complementary raw-patch complexity diagnostic, not the same color scale or phenomenon as local dimension or fiber irregularity.}
\end{figure}
\clearpage

\begin{figure}[p]
\centering
\widepaperfig{\figroot/coco_var_generation_polysemy_aggregate_bins}
\caption{Aggregate VAR-d30 generation diagnostics across all 4096 teacher-forced visual tokens. Left/middle: next-token entropy and NLL binned by local dimension decile. Right: entropy and NLL grouped by fiber-violation strength. The main signal is broad dimension variation, not corrected fiber violations alone.}
\label{fig:coco-var-generation-aggregate}
\end{figure}
\clearpage

\begin{figure}[p]
\centering
\widepaperfig{\figroot/coco_var_generation_polysemy_entropy_heatmaps}
\caption{VAR-d30 predictive-entropy heatmaps from W\&B run \texttt{skt3dcv6}. Entropy is projected back to image-token positions. Bright uncertainty means the generator is less confident about the next visual code, not that the token is geometrically irregular.}
\label{fig:coco-var-generation}
\end{figure}
\clearpage

\begin{figure}[p]
\ContinuedFloat
\centering
\widepaperfig{\figroot/coco_var_generation_polysemy_nll_heatmaps}
\caption{VAR-d30 negative-log-likelihood heatmaps. Like entropy, this is a generation uncertainty diagnostic rather than a fiber-irregularity score.}
\end{figure}
\clearpage

\begin{figure}[p]
\ContinuedFloat
\centering
\widepaperfig{\figroot/coco_var_generation_polysemy_patch_gallery}
\caption{VAR-d30 patch-token gallery for the generation probe. These crops make the selected generator tokens visible so the uncertainty and geometry diagnostics can be read against image content.}
\end{figure}
\clearpage

\begin{figure}[p]
\ContinuedFloat
\centering
\widepaperfig{\figroot/coco_var_generation_polysemy_scatter}
\caption{VAR-d30 geometry/generation scatter. The run-level correlations are weak for irregularity versus generation uncertainty: Pearson Corr($\irr$, entropy) $=-0.033$ and Corr($\irr$, NLL) $=-0.020$. Dimension is more informative for this run, with Spearman Corr($\dimhat$, entropy) $=-0.303$ and Corr($\dimhat$, NLL) $=-0.122$.}
\end{figure}
\clearpage

\begin{figure}[p]
\centering
\widepaperfig{\figroot/coco_var_generation_polysemy_branch_samples}
\caption{Branch samples for matched fiber-violating and control tokens from the 16-image VAR-d30 rerun. Across four matched pairs, singular anchors have higher mean pairwise crop diversity than controls (ratio 2.19), but the sign test is not significant with only four pairs ($p=1.0$). Branch-probability entropy is slightly lower for singular anchors on average. This remains an exploratory diagnostic: it suggests that generation branching should be tested over many more anchors and with stronger image-level diversity metrics rather than inferred from one visually interesting token.}
\label{fig:coco-var-branch}
\end{figure}
\clearpage

\section{Discussion}
Dense vision representations let us test stratified geometry at patch resolution. Each token is both a point in representation space and a visible region in an image, which lets us separate facts that a global intrinsic-dimension estimate would collapse.

The core pattern is stable. DINOv3 and SAM have similar median local dimension on \dataset{COCO}, but differ in same-image locality and slope-change structure. SigLIP2 and AIMv2 move the comparison into a higher-dimensional regime while preserving strong locality. Cross-dataset probes show that the diagnostic changes with visual domain: \dataset{DTD} is sparse-simple for several encoders but not SAM, while \dataset{EuroSAT} makes SAM's corrected fiber violations and sparse burden stand out. A generic change point is not the same as a fiber violation, and sparse raw-patch reconstruction is not the same as local dimension. The VAR-d30 probe adds a useful caution: generation uncertainty aligns more with broad local dimension than with corrected fiber violations.

The restrained claim is enough: modern pretrained vision checkpoints exhibit heterogeneous local geometry, so a single global dimension estimate is a poor summary. This is a diagnostic case study of stratified structure in dense vision spaces, not a final benchmark over all vision models or domains.

\section{Limitations}
The study is intentionally small. The main comparison uses 16 image-aligned \dataset{COCO} images and a limited set of large pretrained checkpoints; the \dataset{STL10}, \dataset{DTD}, and \dataset{EuroSAT} runs are stress tests, not benchmark-scale studies; and VAR-d30 is auxiliary rather than a matched encoder baseline. The local dimension estimator depends on neighborhood range, distance metric, smoothing, and multiple-testing correction. The sparse dictionary probe depends on dictionary size, sparsity cap, and residual threshold. SAM masks provide useful context on \dataset{COCO}, but the fiber statistics are computed from encoder tokens rather than mask quality; outside detection datasets, SAM should be read only as a frozen encoder. A stronger version should scale the image count, compare more checkpoints and datasets, calibrate the change-point test with stronger null controls, and compare against alternative intrinsic-dimension estimators \citep{facco2017estimating}.

\section{Conclusion}
Vision is a dense, spatially grounded medium for the stratified-manifold question raised by language-token embeddings. Across \dataset{COCO}, \dataset{STL10}, \dataset{DTD}, and \dataset{EuroSAT}, frozen DINOv3, SAM, SigLIP2, and AIMv2 patch-token spaces are not well summarized by one global intrinsic dimension. The useful story is local: dimension estimates, fiber-violation maps, sparse-complexity heatmaps, patch--token distance diagnostics, and SAM mask previews together reveal stratified structure that a single number would hide.

\begin{ack}
Acknowledgments are omitted for anonymous review.
\end{ack}

\bibliographystyle{plainnat}
\bibliography{refs}

\appendix

\section{Patch--Token Distance Heatmaps}
\label{app:patch-token-distance}

Raw image patches and model tokens cannot be subtracted directly: a $16\times16$ RGB patch is a pixel vector, while a DINOv3, SAM, SigLIP2, or AIMv2 token is a learned feature vector with a different dimension and scale. The question is therefore not whether one raw patch vector equals one token vector. The question is whether the \emph{geometry of nearby patches} is preserved after the model maps patches into token space.

For each image, we compute two pairwise distance matrices over the same patch grid:
\[
  D^x_{ab}=\|x_a-x_b\|_2,
  \qquad
  D^z_{ab}=\|z_a-z_b\|_2,
\]
where $x_a,x_b$ are raw RGB patches and $z_a,z_b$ are the corresponding encoder tokens. Using Euclidean distance in token space matches the geometry used by the fiber probe. An earlier centered-cosine version overemphasized patch texture and is not used in the paper figures. For the paired DINOv3/SAM comparison, we recompute SAM-H tokens on the exact saved DINOv3-H+ image tensor, so both model heatmaps use the same 16 source images in the same order. SigLIP2-B and AIMv2-L use their own image-aligned 16-image embeddings from the vision-expansion runs.

This distance diagnostic is complementary to the fiber and sparse probes. The fiber probe asks whether token-space volume growth follows the tested local geometric null. The sparse probe asks how hard it is to reconstruct raw patches inside a token neighborhood. Patch--token distance agreement asks a simpler alignment question: when raw patches are close or far in pixel space, are their learned tokens close or far in feature space? A representation can be highly image-local but still strongly reorganize distances within that image.

We use three summaries. First, for each anchor patch, we compare the rank order of distances from that anchor to all other patches with the rank order of distances from the corresponding token to all other tokens. The heatmap value is the within-image Spearman correlation
\[
  \rho_i,
\]
between the raw-patch distance row $D^x_{i,:}$ and the token-distance row $D^z_{i,:}$. Brighter cells mean the token preserves the raw-patch distance ordering more closely; darker cells mean the token neighborhood has moved away from raw RGB similarity. This row-wise score is global within one image: it asks whether the full ordering of all other patches around the anchor is similar in raw and token space.

Second, we report top-neighbor overlap. For each anchor, we compare the nearest 16 raw patches with the nearest 16 tokens and plot the fraction of neighbors that agree. This is a more local diagnostic than the rank-correlation score: it ignores the ordering of distant patches and asks whether the immediate neighborhood is preserved.

Third, the matrix-gallery panels compare the whole normalized raw-patch distance matrix with the whole normalized token-distance matrix. The matrix-level rank correlation summarizes whether the image's overall patch-distance layout is preserved. This view is useful because two models can have similar nearest-neighbor overlap but differ in how they arrange medium- and long-range patch relations.

These scores are not model-quality scores. A low value can mean that the representation has intentionally reorganized raw pixels according to semantics, context, segmentation cues, or invariances. A high value means the feature metric retains more raw-patch geometry. Across the shared 16-image \dataset{COCO} set, DINOv3-H+ has mean distance-rank correlation 0.299, top-16 neighbor overlap 0.279, and matrix-level rank correlation 0.304. SAM-H has 0.565, 0.381, and 0.591. SigLIP2-B has 0.136, 0.168, and 0.146; AIMv2-L has 0.099, 0.119, and 0.095. Thus SAM-H preserves more within-image raw-patch distance structure than DINOv3-H+, while SigLIP2-B and AIMv2-L are more strongly warped away from raw RGB patch geometry.

This interpretation also clarifies the relationship to same-image locality in the main tables. Same-image locality asks whether the nearest token neighbors come from the same source image. Patch--token distance agreement asks a stricter question: within that image, does the token metric preserve which raw patches are close to which other raw patches? DINOv3-H+ is extremely image-local but has lower patch--token distance agreement than SAM-H, suggesting that its feature space keeps image identity while reorganizing within-image patch relations more strongly. SAM-H mixes images earlier, but within each image its token distances track raw-patch distances more closely under this diagnostic.

\paragraph{Readable figure design.}
The paper-facing distance visualizations separate the source image from the score grid. This avoids the color/image interference of overlay heatmaps while preserving the patch layout. The generated full overlay heatmaps and matrix galleries remain available in the artifact directories as audit outputs, but the figures below are the intended reader path.

\begin{figure}[p]
\centering
\widepaperfig{\figroot/coco_patch_token_distance_digest}
\caption{Compact summary of raw-patch versus token-distance agreement across the four \dataset{COCO} encoder runs. Higher values mean that learned token distances preserve more of the within-image raw RGB patch-distance structure. The per-image dot plots show that SAM-H has consistently higher raw-geometry retention than DINOv3-H+, while SigLIP2-B and AIMv2-L are more strongly reorganized away from raw RGB patch distances.}
\label{fig:app-patch-token-distance-digest}
\end{figure}

\begin{figure}[p]
\centering
\widepaperfig{\figroot/coco_dino_sam_patch_token_rank_comparison}
\caption{DINOv3-H+ versus SAM-H distance-rank agreement on representative shared \dataset{COCO} images. The source image is shown separately from the score grids; all score panels use one shared color scale, and the number inside each panel is the per-image mean. This view makes the main qualitative result easier to read: SAM-H more often preserves the raw-patch distance ordering within an image, while DINOv3-H+ reorganizes within-image patch relations more strongly.}
\label{fig:app-dino-sam-patch-token-rank-comparison}
\end{figure}

\begin{figure}[p]
\centering
\widepaperfig{\figroot/coco_dino_sam_patch_token_neighbor_comparison}
\caption{DINOv3-H+ versus SAM-H top-16 neighbor-overlap comparison on representative shared \dataset{COCO} images. Each score cell is the fraction of nearest raw RGB patch neighbors that are also nearest token neighbors for the same anchor patch. This local-neighborhood view complements the rank-agreement figure: SAM-H generally recovers more raw nearest neighbors, but both models retain only a partial copy of raw pixel geometry.}
\label{fig:app-dino-sam-patch-token-neighbor-comparison}
\end{figure}

\clearpage


\section{Residual-Threshold Sweep Visualizations}
\label{app:residual-sweep-visualizations}

This appendix gives the complete visual output for the high-capacity fixed-$k=128$ sparse dictionary residual-threshold sweep. The main text uses $\tau=0.30$ because it preserves regional contrast while avoiding the overly permissive look of very loose residual targets. Here we include every threshold used in the DINOv3-H+/SAM-H sweep, together with the detailed probe dashboard for each model and threshold. All residual-sweep models use 128 local PCA dictionary atoms with maximum OMP sparsity 64. We also include the corresponding SigLIP2-B/AIMv2-L sweep summaries and their $\tau=0.30$ paper-facing heatmaps.

\begin{figure}[H]
\centering
\begin{subfigure}[t]{0.49\linewidth}
  \centering
  \includegraphics[width=\linewidth]{\figroot/coco_dinov3_huge_hicap_residual_sweep}
  \caption{DINOv3-H+}
\end{subfigure}
\hfill
\begin{subfigure}[t]{0.49\linewidth}
  \centering
  \includegraphics[width=\linewidth]{\figroot/coco_sam_hicap_residual_sweep}
  \caption{SAM-H}
\end{subfigure}
\caption{Per-model residual-threshold sweep summaries. Each panel tracks mean and median required sparsity, the 10--90\% band, the fraction of individual patch reconstructions that hit the sparsity cap, and the residual hit rate as the OMP target is relaxed.}
\label{fig:app-residual-sweep-model-summaries}
\end{figure}

\newcommand{\residualheatmapfig}[5]{
\begin{figure}[p]
\centering
\widepaperfig{\figroot/#3}
\caption{#1 sparse-complexity heatmaps at residual target $\tau=#2$. Each panel projects the mean required OMP sparsity for fixed-$k=128$ token neighborhoods back to a $4\times4$ \dataset{COCO} image grid. #4}
\label{#5}
\end{figure}
}

\newcommand{\residualprobefig}[5]{
\begin{figure}[p]
\centering
\includegraphics[width=\linewidth,height=0.78\textheight,keepaspectratio]{\figroot/#3}
\caption{#1 detailed sparse-dictionary probe dashboard at residual target $\tau=#2$. The dashboard summarizes the token-level required sparsity distribution, residual-threshold hit behavior, correlations with local geometry diagnostics, and representative high-complexity patch neighborhoods. #4}
\label{#5}
\end{figure}
}

\subsection{Residual Target tau=0.15}
\residualheatmapfig{DINOv3-H+}{0.15}{coco_dinov3_huge_hicap_tau_0p15_sparse_dictionary_heatmaps}{This strict target remains mostly below the 64-atom cap, but many regions still need several dozen local atoms.}{fig:app-dino-tau015-heatmap}
\residualheatmapfig{SAM-H}{0.15}{coco_sam_hicap_tau_0p15_sparse_dictionary_heatmaps}{This strict target is now mostly free of cap saturation, showing that SAM neighborhoods remain genuinely demanding even with a larger dictionary.}{fig:app-sam-tau015-heatmap}
\residualprobefig{DINOv3-H+}{0.15}{coco_dinov3_huge_hicap_tau_0p15_sparse_dictionary_probe}{The strict target is useful as a stress test for local raw-patch reconstruction.}{fig:app-dino-tau015-probe}
\residualprobefig{SAM-H}{0.15}{coco_sam_hicap_tau_0p15_sparse_dictionary_probe}{The cap-hit statistics show that the high-capacity setting removes most artificial saturation.}{fig:app-sam-tau015-probe}

\clearpage
\subsection{Residual Target tau=0.20}
\residualheatmapfig{DINOv3-H+}{0.20}{coco_dinov3_huge_hicap_tau_0p20_sparse_dictionary_heatmaps}{The regional pattern remains similar to $\tau=0.15$, while the required atom counts decrease.}{fig:app-dino-tau020-heatmap}
\residualheatmapfig{SAM-H}{0.20}{coco_sam_hicap_tau_0p20_sparse_dictionary_heatmaps}{SAM still needs many atoms in visually complex regions, but cap hits are rare.}{fig:app-sam-tau020-heatmap}
\residualprobefig{DINOv3-H+}{0.20}{coco_dinov3_huge_hicap_tau_0p20_sparse_dictionary_probe}{The diagnostic panel shows the same geometry/sparsity relationship under a less strict residual target.}{fig:app-dino-tau020-probe}
\residualprobefig{SAM-H}{0.20}{coco_sam_hicap_tau_0p20_sparse_dictionary_probe}{The hit-rate and cap-share curves improve relative to $\tau=0.15$.}{fig:app-sam-tau020-probe}

\clearpage
\subsection{Residual Target tau=0.25}
\residualheatmapfig{DINOv3-H+}{0.25}{coco_dinov3_huge_hicap_tau_0p25_sparse_dictionary_heatmaps}{DINO keeps visible regional variation without cap effects.}{fig:app-dino-tau025-heatmap}
\residualheatmapfig{SAM-H}{0.25}{coco_sam_hicap_tau_0p25_sparse_dictionary_heatmaps}{SAM shows clear regional differences with essentially no cap-limited neighborhoods.}{fig:app-sam-tau025-heatmap}
\residualprobefig{DINOv3-H+}{0.25}{coco_dinov3_huge_hicap_tau_0p25_sparse_dictionary_probe}{The q10--q90 band remains broad enough to distinguish easier and harder neighborhoods.}{fig:app-dino-tau025-probe}
\residualprobefig{SAM-H}{0.25}{coco_sam_hicap_tau_0p25_sparse_dictionary_probe}{The dashboard illustrates the transition from strict-target stress test toward interpretable regional contrast.}{fig:app-sam-tau025-probe}

\clearpage
\subsection{Residual Target tau=0.30}
\residualheatmapfig{DINOv3-H+}{0.30}{coco_dinov3_huge_hicap_tau_0p30_sparse_dictionary_heatmaps}{This is the DINO threshold used for the main sparse-complexity heatmap because it gives readable contrast without cap saturation.}{fig:app-dino-tau030-heatmap}
\residualheatmapfig{SAM-H}{0.30}{coco_sam_hicap_tau_0p30_sparse_dictionary_heatmaps}{This is the SAM threshold used for the main sparse-complexity heatmap because regional differences remain visible with no cap saturation.}{fig:app-sam-tau030-heatmap}
\residualprobefig{DINOv3-H+}{0.30}{coco_dinov3_huge_hicap_tau_0p30_sparse_dictionary_probe}{The dashboard corresponds to the main-paper DINO sparse heatmap.}{fig:app-dino-tau030-probe}
\residualprobefig{SAM-H}{0.30}{coco_sam_hicap_tau_0p30_sparse_dictionary_probe}{The dashboard corresponds to the main-paper SAM sparse heatmap.}{fig:app-sam-tau030-probe}

\clearpage
\subsection{Residual Target tau=0.35}
\residualheatmapfig{DINOv3-H+}{0.35}{coco_dinov3_huge_hicap_tau_0p35_sparse_dictionary_heatmaps}{Relaxing the target further lowers the required atom counts, while the easiest and hardest regions remain visually distinguishable.}{fig:app-dino-tau035-heatmap}
\residualheatmapfig{SAM-H}{0.35}{coco_sam_hicap_tau_0p35_sparse_dictionary_heatmaps}{SAM has no cap saturation at this target, so the map primarily reflects relative local reconstruction burden.}{fig:app-sam-tau035-heatmap}
\residualprobefig{DINOv3-H+}{0.35}{coco_dinov3_huge_hicap_tau_0p35_sparse_dictionary_probe}{The diagnostic panel shows how the sparsity distribution continues to shift downward.}{fig:app-dino-tau035-probe}
\residualprobefig{SAM-H}{0.35}{coco_sam_hicap_tau_0p35_sparse_dictionary_probe}{The diagnostic panel shows cap-free sparse-complexity contrast.}{fig:app-sam-tau035-probe}

\clearpage
\subsection{Residual Target tau=0.40}
\residualheatmapfig{DINOv3-H+}{0.40}{coco_dinov3_huge_hicap_tau_0p40_sparse_dictionary_heatmaps}{This most relaxed target gives the lowest sparsity counts in the sweep and serves as a high-tolerance comparison.}{fig:app-dino-tau040-heatmap}
\residualheatmapfig{SAM-H}{0.40}{coco_sam_hicap_tau_0p40_sparse_dictionary_heatmaps}{SAM is cap-free at this target, and remaining contrast reflects coarser local reconstruction difficulty.}{fig:app-sam-tau040-heatmap}
\residualprobefig{DINOv3-H+}{0.40}{coco_dinov3_huge_hicap_tau_0p40_sparse_dictionary_probe}{The dashboard completes the DINO residual-threshold sweep.}{fig:app-dino-tau040-probe}
\residualprobefig{SAM-H}{0.40}{coco_sam_hicap_tau_0p40_sparse_dictionary_probe}{The dashboard completes the SAM residual-threshold sweep.}{fig:app-sam-tau040-probe}

\clearpage

\subsection{Additional Vision Encoders}
\begin{figure}[H]
\centering
\widepaperfig{\figroot/coco_vision_expansion_residual_sweep_comparison}
\caption{SigLIP2-B and AIMv2-L residual-threshold sweep comparison. Both use fixed $k=128$ neighborhoods, 128 local PCA atoms, OMP cap 64, and the same threshold grid as the DINOv3/SAM sweep.}
\label{fig:app-vision-expansion-residual-sweep}
\end{figure}

\begin{figure}[p]
\centering
\widepaperfig{\figroot/coco_siglip2_base_tau_0p30_sparse_dictionary_heatmaps}
\caption{SigLIP2-B sparse-complexity heatmaps at the main-paper residual target $\tau=0.30$.}
\label{fig:app-siglip-tau030-heatmap}
\end{figure}

\begin{figure}[p]
\centering
\includegraphics[width=\linewidth,height=0.78\textheight,keepaspectratio]{\figroot/coco_siglip2_base_tau_0p30_sparse_dictionary_probe}
\caption{SigLIP2-B detailed sparse-dictionary probe dashboard at $\tau=0.30$.}
\label{fig:app-siglip-tau030-probe}
\end{figure}

\begin{figure}[p]
\centering
\widepaperfig{\figroot/coco_aimv2_large_tau_0p30_sparse_dictionary_heatmaps}
\caption{AIMv2-L sparse-complexity heatmaps at the main-paper residual target $\tau=0.30$.}
\label{fig:app-aimv2-tau030-heatmap}
\end{figure}

\begin{figure}[p]
\centering
\includegraphics[width=\linewidth,height=0.78\textheight,keepaspectratio]{\figroot/coco_aimv2_large_tau_0p30_sparse_dictionary_probe}
\caption{AIMv2-L detailed sparse-dictionary probe dashboard at $\tau=0.30$.}
\label{fig:app-aimv2-tau030-probe}
\end{figure}

\clearpage


\section{Additional Vision Encoder Geometry Visualizations}
\label{app:additional-vision-encoder-visualizations}

\begin{figure}[p]
\centering
\widepaperfig{\figroot/coco_siglip2_base_fiber_dimension_heatmap}
\caption{SigLIP2-B local dimension heatmap. Color marks the first-regime log-volume/log-radius slope for each patch token.}
\label{fig:app-siglip-dimension-heatmap}
\end{figure}

\begin{figure}[p]
\centering
\widepaperfig{\figroot/coco_aimv2_large_fiber_dimension_heatmap}
\caption{AIMv2-L local dimension heatmap. AIMv2 uses a $16\times16$ analysis grid, yielding 4096 tokens across the same 16-image subset.}
\label{fig:app-aimv2-dimension-heatmap}
\end{figure}

\begin{figure}[p]
\centering
\widepaperfig{\figroot/coco_siglip2_base_dimension_patch_gallery}
\caption{SigLIP2-B patch crops selected by low, middle, high local dimension and high corrected slope-increase score.}
\label{fig:app-siglip-dimension-gallery}
\end{figure}

\begin{figure}[p]
\centering
\widepaperfig{\figroot/coco_aimv2_large_dimension_patch_gallery}
\caption{AIMv2-L patch crops selected by low, middle, high local dimension and high corrected slope-increase score.}
\label{fig:app-aimv2-dimension-gallery}
\end{figure}

\begin{figure}[p]
\centering
\widepaperfig{\figroot/coco_siglip2_base_fiber_volume_curve_gallery}
\caption{Representative SigLIP2-B fiber-bundle volume curves projected back to source image patches.}
\label{fig:app-siglip-volume-curves}
\end{figure}

\begin{figure}[p]
\centering
\widepaperfig{\figroot/coco_aimv2_large_fiber_volume_curve_gallery}
\caption{Representative AIMv2-L fiber-bundle volume curves projected back to source image patches.}
\label{fig:app-aimv2-volume-curves}
\end{figure}

\begin{figure}[p]
\centering
\includegraphics[width=\linewidth]{\figroot/coco_siglip2_base_sparse_dictionary_probe}
\caption{SigLIP2-B sparse-probe dashboard at $\tau=0.30$.}
\label{fig:app-siglip-sparse-dashboard}
\end{figure}

\begin{figure}[p]
\centering
\includegraphics[width=\linewidth]{\figroot/coco_aimv2_large_sparse_dictionary_probe}
\caption{AIMv2-L sparse-probe dashboard at $\tau=0.30$.}
\label{fig:app-aimv2-sparse-dashboard}
\end{figure}

\clearpage

\section{Cross-Dataset Stress-Test Visualizations}
\label{app:cross-dataset-visualizations}

\begin{figure}[p]
\centering
\widepaperfig{\figroot/cross_stl10_dinov3_huge_fiber_irregularity_heatmap}
\caption{\dataset{STL10} DINOv3-H+ corrected fiber-violation heatmap. DINOv3 remains almost perfectly image-local on this object-centered dataset, while the projected violations still appear in spatially localized patches.}
\label{fig:app-cross-stl10-dino-irregularity}
\end{figure}

\begin{figure}[p]
\centering
\widepaperfig{\figroot/cross_stl10_sam_fiber_irregularity_heatmap}
\caption{\dataset{STL10} SAM-H corrected fiber-violation heatmap. SAM has a higher \dataset{STL10} fiber-violation ratio than DINOv3-H+, but these non-detection images do not provide box-prompt mask context, so the plot should be read as encoder geometry rather than segmentation behavior.}
\label{fig:app-cross-stl10-sam-irregularity}
\end{figure}

\begin{figure}[p]
\centering
\widepaperfig{\figroot/cross_dtd_dinov3_huge_sparse_dictionary_heatmaps}
\caption{\dataset{DTD} DINOv3-H+ sparse-complexity heatmap. Texture neighborhoods are relatively easy for the local raw-patch dictionary under this protocol, with mean $S_i=6.29$.}
\label{fig:app-cross-dtd-dino-sparse}
\end{figure}

\begin{figure}[p]
\centering
\widepaperfig{\figroot/cross_dtd_sam_sparse_dictionary_heatmaps}
\caption{\dataset{DTD} SAM-H sparse-complexity heatmap. Compared with DINOv3-H+, SAM-H has much larger raw-patch reconstruction burden on texture images, with mean $S_i=22.92$ under the same local dictionary probe.}
\label{fig:app-cross-dtd-sam-sparse}
\end{figure}

\begin{figure}[p]
\centering
\widepaperfig{\figroot/cross_eurosat_sam_fiber_irregularity_heatmap}
\caption{\dataset{EuroSAT} SAM-H corrected fiber-violation heatmap. This is the strongest cross-dataset stress case in the current runs: SAM-H has low median local dimension but high generic change rate and high corrected slope-increase violation ratio.}
\label{fig:app-cross-eurosat-sam-irregularity}
\end{figure}

\begin{figure}[p]
\centering
\widepaperfig{\figroot/cross_eurosat_dinov3_huge_sparse_dictionary_heatmaps}
\caption{\dataset{EuroSAT} DINOv3-H+ sparse-complexity heatmap. DINOv3-H+ stays strongly image-local on satellite scenes and has a moderate-to-high sparse reconstruction burden, distinct from the lower SigLIP2/AIMv2 sparse scores on the same dataset.}
\label{fig:app-cross-eurosat-dino-sparse}
\end{figure}

\begin{figure}[p]
\centering
\widepaperfig{\figroot/cross_eurosat_sam_sparse_dictionary_heatmaps}
\caption{\dataset{EuroSAT} SAM-H sparse-complexity heatmap. SAM-H has the largest cross-dataset sparse burden in \Cref{tab:cross-dataset-stress}, suggesting that satellite patches are hard for its local raw-patch dictionaries even when the estimated median feature-space dimension is low.}
\label{fig:app-cross-eurosat-sam-sparse}
\end{figure}

\clearpage

\section{Artifact and Reproducibility Statement}
\begin{itemize}[leftmargin=1.5em]
\item The main \dataset{COCO} DINOv3-H+ figures are vendored from \url{runs/local/coco_dinov3_huge_sparse_fiber/20260509_190534}; the W\&B run ID is \texttt{ingbtg9a}.
\item The main \dataset{COCO} SAM ViT-Huge figures are vendored from \url{runs/local/coco_sam_fiber/20260509_191042}; the W\&B run ID is \texttt{nsbl2nhz}.
\item The SigLIP2-B figures are vendored from \url{runs/local/coco_siglip2_base_sparse_fiber/20260509_203403}; the W\&B run ID is \texttt{rqsw6hyt}.
\item The AIMv2-L figures are vendored from \url{runs/local/coco_aimv2_large_sparse_fiber/20260509_203836}; the W\&B run ID is \texttt{xqnbkzlm}.
\item The patch--token distance diagnostics for DINOv3-H+/SAM-H are generated from a shared DINOv3-H+ image tensor. SigLIP2-B and AIMv2-L distance diagnostics use their corresponding image-aligned 16-image embeddings; raw outputs are saved under \url{runs/local/patch_token_distance/20260509_v3} and \url{runs/local/patch_token_distance/20260509_vision_expansion}. The compact digest and clean DINO/SAM comparison grids are generated by \url{scripts/patch_token_distance_digest.py} and \url{scripts/patch_token_distance_comparison_grid.py}.
\item The high-capacity residual-threshold sparse-coding sweep comes from \url{runs/local/sparse_residual_sweep/20260509_hicap}; it reuses the saved DINOv3-H+ and SAM-H embeddings from the two runs above, uses fixed $k=128$ neighborhoods with 128 local PCA atoms and OMP cap 64, and varies only the OMP residual target.
\item The SigLIP2-B/AIMv2-L high-capacity residual-threshold sparse-coding sweep comes from \url{runs/local/sparse_residual_sweep/20260509_vision_expansion}.
\item The SAM ViT-Huge patch-size sensitivity table comes from \url{runs/local/coco_sam_patch_size_sweep/20260508}.
\item The auxiliary \dataset{COCO} VAR-d30 geometry and generation-branching figures are vendored from \url{runs/local/coco_var_d30_sparse_fiber/20260509_194508}; the W\&B run ID is \texttt{3zmdg9wf}. The earlier generation-diagnostic run \texttt{skt3dcv6} remains referenced for continuity.
\item The DINOv3 ViT-L/16 comparison row comes from \url{runs/local/coco_pretrained_vit_sparse_fiber/20260509_193911}; the W\&B run ID is \texttt{dzwnvdkk}.
\item The cross-dataset summary table is saved at \url{runs/local/cross_dataset_logs/cross_dataset_summary_20260511.csv}. The \dataset{STL10} runs are DINOv3-H+ \url{runs/local/stl10_dinov3_huge_sparse_fiber/20260511_102511} (W\&B \texttt{cw9js1pg}), SAM-H \url{runs/local/stl10_sam_huge_sparse_fiber/20260511_102912} (W\&B \texttt{3br3f0da}), SigLIP2-B \url{runs/local/stl10_siglip2_base_sparse_fiber/20260511_101846} (W\&B \texttt{9fld8bxk}), and AIMv2-L \url{runs/local/stl10_aimv2_large_sparse_fiber/20260511_103310} (W\&B \texttt{xf12rlif}).
\item The \dataset{DTD} cross-dataset runs are DINOv3-H+ \url{runs/local/dtd_dinov3_huge_sparse_fiber/20260511_103803} (W\&B \texttt{2z0r30qt}), SAM-H \url{runs/local/dtd_sam_huge_sparse_fiber/20260511_104159} (W\&B \texttt{afs0rz4h}), SigLIP2-B \url{runs/local/dtd_siglip2_base_sparse_fiber/20260511_104554} (W\&B \texttt{1dxplxyq}), and AIMv2-L \url{runs/local/dtd_aimv2_large_sparse_fiber/20260511_104951} (W\&B \texttt{6q5z8po8}).
\item The \dataset{EuroSAT} cross-dataset runs are DINOv3-H+ \url{runs/local/eurosat_dinov3_huge_sparse_fiber/20260511_105415} (W\&B \texttt{5opco63m}), SAM-H \url{runs/local/eurosat_sam_huge_sparse_fiber/20260511_105808} (W\&B \texttt{ov96uga3}), SigLIP2-B \url{runs/local/eurosat_siglip2_base_sparse_fiber/20260511_110200} (W\&B \texttt{mfung05h}), and AIMv2-L \url{runs/local/eurosat_aimv2_large_sparse_fiber/20260511_110549} (W\&B \texttt{bkmujblm}).
\item The final public release should add the exact environment file, model-loading commands, run logs, and the full aggregated JSON/markdown reports for all \dataset{COCO} and cross-dataset runs.
\end{itemize}

\input{checklist.tex}

\end{document}
